% ============================================================
%  Subfile: Pembentukan Ontologi (pilot — kandidat subbab tersendiri)
%  Isi: metodologi pengembangan ontologi, skema (via \subfile{onthology}),
%       dan justifikasi pemilihan OWL.
%  Dipakai oleh bab-3.tex via \subfile{pembentukan-ontologi}, dan dapat
%  dikompilasi mandiri menjadi pembentukan-ontologi.pdf
%  (pdflatex -> biber -> pdflatex -> pdflatex, cwd = folder ini).
% ============================================================
\documentclass[../../thesis]{subfiles}
\begin{document}

% Standalone-pilot cosmetics only; skipped entirely when \subfile'd into thesis.
% Without a chapter context, report-class would label the subsection "0.0.1";
% drop section numbering for the standalone PDF so headings read cleanly.
\ifSubfilesClassLoaded{\setcounter{secnumdepth}{0}}{}

\subsection{Skema Ontologi Knowledge Graph}
\label{subsec:skema-ontologi}

Ontologi dalam penelitian ini diposisikan menggunakan dua tier
representasi yang dibedakan secara metodologis. \textbf{Tier
ekstraksi} menghasilkan model konseptual \textit{deskriptif} atas isi
buku teks Kurikulum Merdeka (konsep, sub-bab, bab, beserta
keterkaitan semantik antar-konsep dalam satu buku), sejalan dengan
tradisi \textit{domain module} dalam konstruksi otomatis \textit{knowledge
graph} dari buku teks elektronik. \textbf{Tier completion} menghasilkan kandidat
keterkaitan lintas-buku (\texttt{LINTAS\_BUKU\_*}) sebagai inferensi
pedagogis dangkal di atas Tier ekstraksi: sistem mengusulkan jembatan
antar-disiplin yang dapat dipertimbangkan, tetapi tidak menetapkannya
sebagai jalur belajar.

Predikat ``dangkal'' bersifat operasional: sistem hanya menginferensi
keterkaitan \textit{single-hop} antar dua konsep di buku berbeda dan
tidak menyusun jalur belajar multi-langkah lintas-disiplin (batasan
inferensi pedagogis dirinci pada Subbab~\ref{sec:ruang-lingkup}).
Dengan demikian, relasi berarah pada Tier ekstraksi seperti
\texttt{PRASYARAT} dan \texttt{MEMPERSIAPKAN} dimaknai sebagai
ketergantungan konseptual yang terdokumentasikan dalam materi sumber
(keputusan editorial penyusun buku), sedangkan relasi lintas-buku pada Tier
completion seperti \texttt{LINTAS\_BUKU\_PRASYARAT\_UNTUK} dimaknai sebagai
sugesti per-pasangan yang dapat dipertimbangkan secara lokal, bukan anjuran
kronologis bagi siswa.

\subsubsection{Metodologi Pengembangan Ontologi}

Ontologi pada penelitian ini tidak dirancang secara manual di muka, melainkan
disusun secara \textbf{induktif} (\textit{bottom-up}) melalui ekstraksi
berbantuan-LLM atas korpus: kelas, hierarki, dan vokabulari relasi diturunkan
dari pemrosesan teks buku, kemudian dikurasi dan divalidasi pakar. Kerangka
kelas hierarkis mengikuti struktur dokumen sumber (daftar isi), sedangkan tipe
relasi muncul dari keluaran ekstraksi lalu dikurasi menjadi himpunan inti.
Sesuai posisinya sebagai ontologi ringan pada spektrum formalitas
(Gambar~\ref{fig:spektrum-ontologi}), ontologi ini cukup formal untuk
menegakkan konsistensi tipe antartahap \textit{pipeline} tanpa memuat aksioma
penalaran berat yang tidak diperlukan oleh tujuan penelitian. Untuk keperluan
pemaparan, proses pengembangan ini \textbf{dapat dipetakan ke} tujuh tahapan
panduan \textit{Ontology Development 101} \parencite{noy_ontology_2001} sebagai
berikut.

\begin{enumerate}
  \item \textbf{Penentuan domain dan lingkup:} konsep sains lintas-disiplin
        (Fisika, Kimia, Biologi) Fase~F Kelas~XII Kurikulum Merdeka.
  \item \textbf{Pertimbangan penggunaan-ulang:} tidak ditemukan ontologi
        kurikulum sains berbahasa Indonesia yang dapat digunakan ulang secara
        langsung, sehingga ontologi dirancang khusus dan sengaja dijaga ringan.
  \item \textbf{Enumerasi istilah penting:} daftar isi Buku Siswa memasok
        struktur bab dan subbab, sedangkan \textit{materi pokok} dari Buku
        Panduan Guru beserta glosarium buku memasok terminologi konsep (rincian
        blok masukan \textit{prompt} pada
        Lampiran~\ref{lamp:prompt-ekstraksi}).
  \item \textbf{Definisi kelas dan hierarki:} rantai
        \texttt{Grade/Subject} $\rightarrow$ \texttt{Chapter} $\rightarrow$
        \texttt{Subtopic} $\rightarrow$ \texttt{Concept} diturunkan dari
        struktur daftar isi; kelas \texttt{ConceptTarget} ditambahkan sebagai
        keputusan rekayasa untuk menampung target relasi yang belum terekstrak
        sebagai \texttt{Concept} lengkap.
  \item \textbf{Definisi properti:} tipe relasi antarkonsep muncul dari hasil
        ekstraksi lalu dikurasi menjadi himpunan inti lima belas tipe
        dalam-buku yang merepresentasikan kategori relasi ilmiah umum (kausal,
        komposisional, fungsional, definisional, dan sebagainya), didampingi
        properti data seperti \texttt{description} (sumber \textit{embedding}),
        \texttt{glossary\_validated}, dan \texttt{materi\_pokok\_ref}.
  \item \textbf{Penetapan batasan properti:} tiap relasi semantik diberi domain
        dan jangkauan \texttt{Concept} serta sifat simetri bila berlaku.
  \item \textbf{Pembuatan instans:} konsep dan relasi aktual diisi oleh
        \textit{pipeline} ekstraksi berbantuan-LLM, kemudian divalidasi enam
        pakar domain dan disempurnakan melalui empat iterasi keadaan KG
        (v1--v4).
\end{enumerate}

Karena konseptualisasi bertumpu pada keluaran ekstraksi dan bukan pada
perancangan manual, panduan \textcite{noy_ontology_2001} di sini berperan
\textbf{deskriptif-analitis} --- sebagai kerangka untuk menstrukturkan
pemaparan proses --- bukan prosedur yang diikuti langkah demi langkah.

Skema \textit{knowledge graph} diformalkan sebagai ontologi OWL untuk
memastikan konsistensi struktur dan vokabulari relasi antartahap
\textit{pipeline} (ekstraksi, konstruksi, dan \textit{completion}).
Topologi keseluruhan mencakup relasi struktural, relasi keterurutan
penyajian materi dalam buku (\textit{NEXT\_CHAPTER}), dan relasi semantik
pada Tier ekstraksi, serta inferensi lintas-buku pada Tier completion.
Keseluruhan topologi ini diilustrasikan pada Gambar~\ref{fig:ontologi}.

\subfile{onthology}

\subsubsection{Justifikasi Pemilihan OWL}

Skema diformalkan menggunakan \textit{Web Ontology Language} (OWL)
\parencite{w3c_owl2_2012}. Pemilihan OWL didasarkan pada apa yang
benar-benar dimanfaatkan darinya. \emph{Pertama}, tipisasi formal yang dapat
diperiksa mesin: deklarasi kelas (\texttt{owl:Class}), properti objek
(\texttt{owl:ObjectProperty}), dan properti data
(\texttt{owl:DatatypeProperty}) beserta batasan
\texttt{rdfs:domain}/\texttt{rdfs:range}, sehingga vokabulari struktur dan
relasi menjadi spesifikasi eksplisit yang konsisten antartahap
\textit{pipeline}. \emph{Kedua}, aksioma semantik ringan: sifat simetri relasi
dinyatakan melalui \texttt{owl:SymmetricProperty} (mis.\
\texttt{BERINTERAKSI\_DENGAN}, \texttt{BEREAKSI\_DENGAN}, dan
\texttt{LINTAS\_BUKU\_SAMA\_DENGAN}). \emph{Ketiga}, sebagai standar terbuka
W3C, OWL didukung perkakas baku seperti Prot\'eg\'e, bersifat portabel, dan
dapat divalidasi serta digunakan ulang oleh peneliti lain.

Alternatif lain dinilai kurang sesuai: RDFS murni tidak menyediakan aksioma
properti seperti simetri; SHACL, JSON Schema, maupun Pydantic berorientasi pada
validasi \textit{bentuk} (\textit{shape}) data, bukan ontologi formal yang
menyatakan semantik kelas dan relasi secara deklaratif; sedangkan UML atau
diagram informal tidak baku secara semantik. Sejalan dengan pembahasan spektrum
formalitas pada Bab~2 (Gambar~\ref{fig:spektrum-ontologi}), menuliskan ontologi
ringan ke dalam OWL adalah hal yang wajar karena OWL menyediakan profil dari
yang ringan hingga berat \parencite{w3c_owl2_profiles_2012}.

OWL berperan sebagai \textbf{spesifikasi waktu-desain}, sedangkan penyimpanan
dan kueri \textit{runtime} memakai \textit{Labeled Property Graph} (LPG) pada
Neo4j; tidak terdapat \textit{triplestore}, SPARQL, maupun \textit{reasoner}
pada jalur produksi. Pemetaannya langsung: \texttt{owl:Class} menjadi label
node, \texttt{owl:ObjectProperty} menjadi tipe relasi,
\texttt{rdfs:domain}/\texttt{rdfs:range} ditegakkan melalui pola \texttt{MERGE}
dan \textit{constraint} Cypher, dan properti data menjadi properti pada node
maupun \textit{edge}.

Vokabulari relasi diarahkan melalui instruksi \textit{prompt}
(Lampiran~\ref{lamp:prompt-ekstraksi}) dan dikonsolidasikan pada tahap
konsensus; cakupan penegakannya dibahas pada keterbatasan penelitian (Bab~6).

\end{document}
